% commentaire.tex — f174v, vv. 1001–1043 (conversion initiale des travaux antérieurs)
\comm{1007}{πέρ᾽ ἐγὼ 'ξινίγκι πορμός, ἴνα πυλάξι σοι}{$=$ φέρ᾽ ἐγὼ 'ξενέγκω φορμόν, ἵνα φυλάξω σε.}
\comm{1015–135}{Parodie de l'\italique{Andromède}.}{La pièce perdue d'Euripide \italique{Andromède} fut peut-être représentée la même année que l'\italique{Hélène}\nc{C'est ce qu'indiquent les scholies au vers 1012, πιθανῶς· συνδεδίδακται γὰρ τῇ Ἑλένη, ainsi que celle du vers 53 des \italique{Grenouilles}. Le début du prologue d'\italique{Andromède} est repris au vers 1065, «~ὦ νὺξ ἱερὰ~», -- oui : Aristophane a chamboulé l'ordre de la pièce ; scholie à ce vers 1065 : τοῦ προλόγου Ἀνδρομέδας εἰσβολή --, laisse à penser qu'\italique{Andromède} fut jouée au lever du jour, avant donc \italique{Hélène}. Matthew Whright (\italique{Euripides' Escape-Tragedies}, 2005) fait l'hypothèse contraire et avance qu'\italique{Iphigénie en Tauride} concluait la trilogie d'Euripide de 412, agrémenté du drame satyrique du \italique{Cyclope}. Voir la recension que feu Alain Blanchard en fit dans la \italique{Revue des Études Grecques}, tome 119 (2006), pp. 798-800.}, en 412, un avant les \italique{Thesmophories}. L'édition des fragments d'Euripide de Jouan \& van Looy\nc{Tome VIII (fragments), première partie, de l'édition d'Euripide aux Belles Lettres (1998).} propose un dossier complet auquel le lecteur pourra se reporter. Le récit le plus circonstancié du mythe d'Andromède se trouve chez Apollodore (II, 4, 3) :\citComm{Παραγενόμενος δὲ εἰς Αἰθιοπίαν, ἧς ἐβασίλευε Κηφεύς, εὗρε τὴν τούτου θυγατέρα Ἀνδρομέδαν παρακειμένην βορὰν θαλασσίῳ κήτει. Κασσιέπεια γὰρ ἡ Κηφέως γυνὴ Νηρηίσιν ἤρισε περὶ κάλλους, καὶ πασῶν εἶναι κρείσσων ηὔχησεν· ὅθεν αἱ Νηρηίδες ἐμήνισαν, καὶ Ποσειδῶν αὐταῖς συνοργισθεὶς πλήμμυράν τε ἐπὶ τὴν χώραν ἔπεμψε καὶ κῆτος. Ἄμμωνος δὲ χρήσαντος τὴν ἀπαλλαγὴν τῆς συμφορᾶς, ἐὰν ἡ Κασσιεπείας θυγάτηρ Ἀνδρομέδα προτεθῇ τῷ κήτει βορά, τοῦτο ἀναγκασθεὶς ὁ Κηφεὺς ὑπὸ τῶν Αἰθιόπων ἔπραξε, καὶ προσέδησε τὴν θυγατέρα πέτρα. Ταύτην θεασάμενος ὁ Περσεὺς καὶ ἐρασθεὶς ἀναιρήσειν ὑπέσχετο Κηφεῖ τὸ κῆτος, εἰ μέλλει σωθεῖσαν αὐτὴν αὐτῷ δώσειν γυναῖκα. Ἐπὶ τούτοις γενομένων ὅρκων, ὑποστὰς τὸ κῆτος ἔκτεινε καὶ τὴν Ἀνδρομέδαν ἔλυσεν. Ἐπιβουλεύοντος δὲ αὐτῷ Φινέως, ὃς ἦν ἀδελφὸς τοῦ Κηφέως ἐγγεγυημένος πρῶτος τὴν Ἀνδρομέδαν, μαθὼν τὴν ἐπιβουλήν, τὴν Γοργόνα δείξας μετὰ τῶν συνεπιβουλευόντων αὐτὸν ἐλίθωσε παραχρῆμα.\\
Parvenu en Éthiopie, dont Céphée était roi, [Persée] trouva sa fille Andromède exposée pour être dévorée par un monstre marin. Cassiopée, épouse de Céphée, avait osé se comparer aux Néréides pour la beauté, et s'était même vantée de l'emporter sur elles. Les Néréides en furent irritées. Neptune partagea leur indignation, submergea le pays, et y envoya un monstre marin. L'oracle d'Ammon ayant annoncé que ces désastres cesseraient si on exposait Andromède, fille de Cassiopée, pour être dévorée par le monstre, les Éthiopiens forcèrent Céphée à faire ce que l'oracle ordonnait et à attacher sa fille à un rocher. Persée l'ayant vu, il devint amoureux, et promit à Céphée de tuer le monstre s'il voulait la lui donner en mariage. Céphée s'y étant engagé par serment, il attendit le monstre, le tua et délivra Andromède. Phinée, frère de Céphée, à qui Andromède avait été promise avant cet événement, conspira contre lui pour le faire périr ; mais Persée l'ayant découvert, lui montra la tête de la Gorgone, et le changea en pierre, ainsi que tous ceux qui avaient pris part à son complot.}}
\comm{1015–55}{Monodie d'Andromède-Mnésiloque.}{La structure de ce chant a fait couler beaucoup d'encre... et nous échappe encore\nc{On pourra aller voir ce qu'en dit Lætitia, pp. 436-445.}. Mais bien loin de l'anarchie décrite par les commentateurs, je suis tenté d'y voir une illustration sonore des échos : chaque phrase musicale semble être répétée : l'on peut ainsi regrouper par paires successives les différentes structures métriques (le thème puis son écho, de mètre identique ou très proche). On regroupe ainsi les vers : 1015-1071 \& 1018-1021 ; 1022-1023 \& 1024-1025 ; 1026-1029 \& 1030-1031 ; 1032-1033 \& 1034-1035 ; 1036-1037 \& 1038-1039a ; 1043 \& 1045 ; 1050-1051 \& 1052-1053.}
\comm{1018}{ὦ προσᾳδοῦσσαι τὰς ἐν ἄντροις.}{Comme dans la pièce d'Euripide, la nymphe Écho répond aux plaintes d'Andromède. (Dans la reprise comique, c'est Euripide qui se métamorphose en Écho.)
Willems a consacré sa \nieme{25} note (reproduite en annexe) à ce vers et aux diverses retouches qu'il a subies.
Je reprends la lecture de Burges\nc{G. Burges, \italique{Classical Journal} 14 (1816) 217. Ou bien voir W. Mitsdörffer, «~Das Mnesichoslied in Aristophanes' \italique{Thesmophoriazusen}~», \italique{Philologus} 98 (1954-55), p. 69.}, προσᾴδουσ' ἀϋτάς, satisfaisante quant au sens et qui offre l'avantage de rester très proche du texte du manuscrit.}
\comm{1026–8}{ὅδε γὰρ ὁ Σκύθης φύλαξ / πάλαι ἐφέστηκ᾽ ὀλοὸν / ἄφιλον ἐκρέμασε κόραξι δεῖπνον.}{Ce passage est sans doute légèrement corrompu. Les tentatives pour le corriger ont toutes été infructueuses\nc{Cf. l'inventaire sans concession qu'en fait Lætitia \italique{The Songs of Aristophanes}, p. 443-444 : son regard infaillible de métricienne aguerrie détecte les invraisemblances métriques des «~corrections~». Sa proposition n'en est pas moins aporétique : elle voit dans ces vers une suite de trois dochmiaques se terminant par un dimètre iambique catalectique tout en concédant que la correption πάλα\lettreBreve{ι} (\latin{vocalis ante vocalem corripitur}) est suspecte : les seuls autres exemples chez Aristophane étant au vers 915, «~κύσ\lettreBreve{ω}. ἄπαγε~», déjà évoqué plus haut, et le peu convaincant «~ἰὼ ἰώ~» au vers 1170 des \italique{Nuées}.}.}
\comm{1031}{ψῆφον κημὸν ἕστηκ᾽ ἔχουσ᾽.}{On mettra ce vers en regard d'un fragment d'Hermippe :\citComm{καὶ πρὸς κύβους ἕστηκ' ἔχων <τὸ> κήθιον.}τὸ κήθιον et ὁ κημός sont synonymes : ils désignent dans l'univers ludique le \italique{cornet} dans lequel on secouait les dés\nc{Cf. Pollux (X, 150) sur les ustensiles de jeu : \citComm{Κυβευτοῦ σκεύη · ἄβαξ ἀβάκιον, κόσκινον, κύβοι, διάσειστοι κύβοι, ἀστράγαλοι, φιμοί, κημοί, τηλία, κήθια κηθίδια, ψῆφοι, πεττοί.}φιμός et κημός sont littéralement des \italique{muselières}, avec le sens spécialisé de \italique{cornets} (à dés ou astragales) et sont synonymes de κήθιον. On pourra par exemple se reporter à l'article de Michel Casevitz, «~Sur les termes de jeu : compléments d'histoire des mots~», \italique{Kentron} 36 (2021), p. 209-212. Pour des ouvrages beaucoup plus conséquents, le lecteur pourra lire \italique{Les jeux des Anciens}\expo{2} (1873) de Louis Becq de Fouquières, auteur à la vaste érudition et au style des plus agréables comme le \nieme{XIX} siècle en a été prodigue ; l'édition par Jean Taillardat déjà mentionnée du Περὶ βλασφημιῶν, \italique{Des termes injurieux} et du Περὶ παιδιῶν, \italique{Des jeux grecs}, de Suétone ; la somme de référence qu'est l'article de Hans Lamer «~\latin{Lusoria tabula}~» de la \italique{Paulys Realencyclopädie der classischen Altertumswissenschaft}, \allemand{Band} XIII 2 (1927), col. 1900-2029 ; plus récemment, les \latin{excerpta de ludis} chez Pollux ont été publiés par Salvatore Costanza : \italique{Giulio polluce, onomasticon: excerpta de ludis} accompagnés d'une traduction italienne. Outre l'archéologie, nos témoins principaux quant aux jeux grecs sont Pollux, Hésychius et Eustathe (qui s'appuie sur les traités perdus de Suétone) -- non exempts de contradictions entre eux d'ailleurs.}. ἡ ψῆφος, comme λίθος, est dans ce contexte un \italique{pion}.
On constate avec ces deux exemples que `ἕστηκεν ἔχων' devait être l'expression consacrée pour le joueur qui s'apprêtait à lancer les dés.
On a soupçonné une altération du vers du fait de l'étrange succession initiale de trois longues. 
L'exemple d'Hermippe suscite naturellement l'émendation suivante :\citComm{ὁρᾷς; οὐ χοροῖσιν οὐδ᾽ ὑφ᾽ ἡλίκων ν\synizese{εα}νίδων\\
ψῆφον ἕστηκ᾽ ἔχουσα κημόν.}(κημός placé après ἔχω, les mots sont ainsi agencés dans le même ordre que dans le fragment\nc{Van Leeuwen a recensé d'autres erreurs similaires où l'ordre des mots est inversé dans le \italique{Ravennas} (la comparaison avec les neuf pièces du \italique{Ravennas} qui nous sont connues par d'autres témoins est préférable : la correction est davantage sûre). Cf. l'édition en fac-similé qu'il fit de notre manuscrit : \italique{Comoediae Undecim cum Scholiis : Codex Ravennas 137, 4, A} [$=$ \italique{Rav.} 429 nouvelle cote], \italique{phototypice editus} in \italique{Codices Graeci et Latini photographice depicti}, t. IX (1904), pp. XVIII-XIX.}.)\\
L'avantage est qu'on obtient (avec la synizèze de ν\synizese{εα}νίδων) le schéma métrique suivant :\citComm{$ \cup \ - \ - \ \cup \ - \ | \ \cup \ - \ \cup \ - \ | \ \cup \ - \ - \ \cup \ -$\\
$ - \ \cup \ - \ | \ - \ \cup \ - \ | \ \cup \ - \ \smile $\\
dochmiaque $+$ iambe $+$ dochmiaque \\
crétique $+$ crétique $+$ bacchée,}de scansion identique à celle des deux vers précédents, annulant du même coup le surprenant abrègement de πάλα\lettreBreve{ι} en seconde position d'une résolution (qui le faisait scander $ \widearc{\cup \cup}$, cf. commentaire v. 1026-8 n.1). 
Je ne sais ce que vaut cette émendation ; je laisse le soin d'en juger à des métriciens compétents.\\
Il est en tout cas bien évident qu'Andromède-Mnésiloque, se lamentant de son sort présent et regrettant les douces années de son enfance, ne fait aucunement allusion à une quelconque \italique{urne de vote} (κημός ?). Les chœurs de danse et les jeux enfantins, voilà qui suffit à sa mélancolie. Le marigot électoral ? alors que la jeune fille ne pouvait même pas y participer activement ? Mais où donc l'ensemble des commentateurs -- à commencer par le scholiaste -- est-il allé cherché cette abracadabrantesque explication ?
Qu'Aristophane ait pu utiliser le mot κημός, nom d'une sorte d'entonnoir posé sur l'urne pour en celer le vote, comme synecdoque pour l'\italique{urne} elle-même, cela est vrai mais avec une réserve de taille ayant échappé aux commentateurs : c'est pour permettre un jeu de mots, comme nous l'allons voir. (Ce qui explique qu'il soit le seul auteur recourant à ce vocable avec le sens d'\italique{urne}.)\\
Cette inadmissible traduction de κημός par «~urne~» [pour le vote] nous oblige à rebondir sur un passage des \italique{Guêpes}. Il s'agit du début des \italique{Guêpes} avec la traditionnelle scène des deux esclaves qui présentent la situation : en l'occurrence, le mal dont souffre leur maître Philocléon (et une bonne partie des Athéniens), pris qu'il est de \italique{philéliastite}. Et lorsqu'il voit une dédicace au beau Dêmos, le fils de Pyrilampe, notre malade ne peut s'empêcher d'écrire en marge (παρέγραψε) : κημὸς καλός. Vers 97-99 :
\citComm{καὶ νὴ Δί᾽ ἢν ἴδῃ γέ που γεγραμμένον\\
υἱὸν Πυριλάμπους ἐν θύρᾳ Δῆμον καλόν,\\
ἰὼν παρέγραψε πλησίον `κημὸς καλός.'}Qui était ce Dêmos Πυριλάμπους ? On possède quelques autres témoignages à son sujet : Calliclès, l'horrible démagogue du \italique{Gorgias} (481d), en avait fait son éromène. 
Platon souligne malicieusement que, à l'instar de l'autre dêmos que lorgne le démagogue, Calliclès finissait toujours par céder devant son cadet aimé (αἰσθάνομαι οὖν σου ἑκάστοτε, καίπερ ὄντος δεινοῦ, ὅτι ἂν φῇ σου τὰ παιδικὰ καὶ ὅπως ἂν φῇ ἔχειν, οὐ δυναμένου ἀντιλέγειν). 
Eupolis (fr. 227(213) Kassel-Austin) nous le décrit comme \italique{bouché} ; littéralement : il est plein de cérumen (κυψέλη) : 
\citComm{καὶ τῷ Πυριλάμπους ἆρα Δήμῳ κυψέλη ἔνεστιν; }Enfin, Athénée (IX, 397c-d), après avoir décrit le paon comme un animal bruyant qui réveille les dormeurs (τοιόνδε ταὧν, ὃς τοὺς εὕδοντας ἐγείρει) nous apprend que le fils de Pyrilampe en faisait élevage, dont il était d'ailleurs très fier...\\
Que peut-on en conclure ? assurément, que le beau Dêmos était un invétéré babillard, obtus à tout conseil, qui en remontrait à Calliclès au rebours d'une relation maître-élève classique et qui, pas plus que les paons qu'il élevait, ne s'arrêtait jamais de parler.
Comment lire maintenant le graffiti : \citComm{κημὸς καλός Δῆμον καλόν}(noter l'accusatif de Δῆμον καλόν) ?
Tout simplement en conformité avec le sens premier de κημός, savoir : une \italique{muselière}.\citComm{Une bonne muselière [fera taire] le beau Dêmos.}Que la muselière (κημός), ait pu désigner aussi une \italique{urne de vote}, voilà qui devait laisser Aristophane bien rêveur !..
La pique contre la démocratie est évidente :
\citComm{Une bonne muselière [musellera] le bon peuple [pour cinq ans, pourrait-on ajouter à propos de la nôtre, encore plus décadente que celle d'Athènes].}Mais Philocléon dans sa naïveté -- qui serait touchante si elle n'était pas criminelle -- ne voit tout ce qu'il y a d'avilissant dans cette mascarade démocratique ; il est satisfait que le bon peuple puisse urner à son aise :
κημὸς καλός δῆμον καλόν.\\
Autre exemple, un peu plus loin dans la même pièce : Philocléon préside au procès du général Lachès\nc{La scène suivante est une amusante parodie du procès intenté, ou censé intenté, au stratège athénien Lachès, qui était en effet du dême d'Æxone. Le chien du peuple, c'est Cléon, du dême de Cydathène. Lachès avait commandé la flotte envoyée en Sicile en 427. Rappelé au commencement de 425, il fut accusé d'avoir malversé. Le procès (s'il eut lieu) fut présidé suivant l'usage par un thesmothète (v. 935). Cléon, comme synégore, soutint l'accusation. Lachès fut acquitté, car il remplit encore diverses fonctions publiques. Il est intéressant de constater que si le poète n'épargne pas Cléon, il n'est guère plus indulgent pour Lachès, encore que celui-ci appartînt au parti aristocratique. (Note de Willems.)}, animalisé pour l'occasion en chien \italique{Chapardeur} (Λάβης). L'accusateur public n'est autre que le chien de Dêmos, Cléon. Bdélycléon fait entrer le prévenu et son accusateur ; c'est-à-dire les deux chiens, la canaille et le roquet. Réaction de Philocléon devant ces deux patibulaires (\italique{Guêpes}, vers 852-853): \citComm{Φιλοκλέων -- ἐς κόρακας, ὡς ἄχθομαι\\
ὁτιὴ 'πελαθόμην τοὺς καδίσκους ἐκφέρειν.\\
Aux corbeaux ! Quel ennui ! j'ai oublié d'apporter les καδίσκοι.}On traduit καδίσκοι par \italique{urnes} ; et il n'est pas douteux que dans la parodie de procès, les καδίσκοι soient mis en place d'\italique{urnes}... mais on voit que Philocléon frémit tout simplement devant les deux molosses qu'on lui présente et regrette l'absence de \italique{muselières}, c'est sans doute aussi le sens propre de καδίσκος.\\}
\comm{1033}{κήτει βορὰ Γλαυκέτῃ πρόκειμαι.}{On sait par le scholiaste des \italique{Oiseaux} (vers 348) que la pièce originale d'Euripide comportait à cet endroit :\citComm{ἐκθεῖναι κήτεϊ φορβάν,}Γλαυκέτης, l'Athénien glouton (τένθης) est cité au vers 1008 de la \italique{Paix}.
Trygée y fait une prière pour que la paix revienne à Athènes ; la ville est sens dessus dessous du fait de la guerre : les embargos et autres razzias ont littéralement renversé le marché : au lieu de toutes les bonnes denrées, l'Athénien n'y rencontre que des rappels de ses privations. 
Autre inversion : si la paix revenait, ce serait alors aux citoyens de bousculer les trois parasites du marché aux anguilles. Aristophane joue sur une image qu'il a développée ailleurs\nc{Cavaliers, vers 864-867. Voir aussi le vers 361. Érasme a relevé l'image dans ses \italique{Adages} : n° 2576 (\latin{Anguillas captare}).} : comme le chasseur d'anguille trouble l'eau pour capturer sa proie, le politicien véreux trouble les affaires publiques pour mieux s'enrichir. (Remarquer l'utilisation du verbe τυρβάζεσθαι, qu'Aristophane emploie aussi au vers 257 des \italique{Guêpes} : τὸν πηλὸν ὥσπερ ἀτταγᾶς τυρβάσεις βαδίζων.) 
On peut mettre notre main au feu que le groupe des trois affreux étaient de francs coquins qui se gobergeaient sur le dos d'Athènes pour mener une vie douce, bonne et grasse\nc{Cf. pour Morykhos les \italique{Acharniens}, v. 887 ainsi que les \italique{Guêpes}, v. 506 \& 1142. On notera au passage qu'un Téléas fils de Télénikos apparaît à la même époque dans les \italique{Inscriptiones Graecae} (I\expo{3}, 307 ; 308 ; 331 ; 332 ; 370) à la secrétairie (ἐγραμμάτευε) du Trésor (ℎοι ταμίαι το͂ν ℎιερο͂ν χρεμάτον τε͂ς Ἀθεναίας). Voir ce qu'en disent les scholies aux vers 167-8 des \italique{Oiseaux}.}.
Γλαυκέτης est mis par paronomase ; \italique{Nomen, omen} : lorsqu'on s'appelle Γλαυκέτης, on est prédestiné à devenir ce poisson vorace γλαῦκος-κῆτος\nc{Autre jeu du même acabit sur le nom de Glaucétès : (Platon le Comique, fr. 114(106) K.-A.) :\citComm{ὦ θεῖε Μόρυχε· πῶς γὰρ οὐ δαίμων ἔφυς;\\
καὶ Γλαυκέτης ἡ ψῆττα, καὶ Λεωγόρας,\\
οἵ ζῆτε τερπνῶς οὐδὲν ἐνθυμούμενοι.}} !
(L'image de la femme livrée en pâture (βορά) à la poissonnaille est récurrente chez les Comiques, \italique{confer} Athénée (VII, 315b) citant Platon le Comique ($=$ fr. 57(56) Kassel-Austin)) :\citComm{σὲ γάρ, γραῦ, συγκατῴκισεν σαπρὰν\\
ὀρφῷσι σελαχίοις τε καὶ φάγροις βοράν}(βορά désigne ici la menue fretaille qui sert de pâture aux poissons carnivores comme l'orphe (ὄρφος / ὀρφώς), qu'Athénée dans le même chapitre rapproche du \italique{glauque}, γλαῦκος.)\\}
\comm{1041}{Ἀίδα,}{Aristophane use du même génitif dorique au vers 1333 des \italique{Grenouilles} (parodie d'Euripide).}
\comm{1041}{πολυδάκρυτον Ἀίδα γόον φεύγουσαν.}{Le manuscrit porte l'incompréhensible †φεύγουσαν†. Y suppléer par χέουσα(ν), à l'instar de certains éditeurs, c'est faire fi de sa leçon. Je dirais même plus : χέουσα est le mot banal qu'on attendait, il paraît alors très difficile qu'un scribe s'y soit trompé. (σφριγῶσα, \italique{gonflée}, aurait au moins le mérite de se justifier paléographiquement mais le terme paraît impropre en grec à l'image d'yeux gonflés de larmes.)
Une autre proposition, `φλέγουσα(ν)', a aussi emporté la faveur de nombreuses éditions. Elle est reprise par Jean Taillardat à propos de l'image de la Muse-flamme (§736 n.2)... mais le grand Taillardat me semble s'abuser : si un poème, un chant, un \italique{thrêne} (qui est avant tout sonore ; la flamme étant celle de la voix selon moi) peuvent être enflammés, cela ne prouve pas pour autant que des pleurs peuvent l'être.\\
Le scholiaste nous fournit une information qui pourrait peut-être bien débrouiller l'écheveau: une glose, σκοτεινήν, \italique{obscur}, se trouve au vers 1054 pour qualifier soi-disant αἰόλαν, \italique{rapide}, \italique{mobile}, \italique{changeant}. 
Sauf qu'αἰόλος n'a jamais eu le sens de σκοτεινός et que Willems a clairement montré qu'il fallait prendre αἰόλος en son sens courant\nc{Note n° 27, reproduite ci-après.}.
Suidas (αι 245) -- ou l'un des lexicographes de la \italique{Souda}, si l'on préfère -- n'a pas manqué néanmoins de relever la glose à l'entrée Αἰόλαν. Mais Suidas est coutumier de ce genre de fausses attributions (voir la note au vers 735).
Or en cherchant des mots proches de †φεύγουσαν†, je suis retombé sur ce cher Hésychius (ψ 135) :\citComm{ψεφαυγοῦς· σκοτεινῆς,}Génitif dorique (comme Ἀίδα) d'un certain *ψεφαυγής, \italique{au sombre regard}.
Euripide a une prédilection pour forger des composés en -αυγής. (\italique{Confer} le Chantraine, \italique{sub voce} αὐγή.) 
Si l'on fait l'hypothèse de la chute de `ψε', le φαυγους restant aurait bien pu être complété d'un `αν' en regard de l'autre participe †λιτομέναν† juste au-dessus, et donner ainsi notre leçon †φεύγουσαν†.}
\comm{1043–51}{Papyrus d'Oxyrhynque n° 4935.}{Un maigre fragment (n° 1) du Papyrus d'Oxyrhynque 4935\nc{\italique{The Oxyrhynchus Papyri}, LXXIII (2009), ed. A. Benaissa.} (\nieme{II} siècle ap. J.-C.), reproduit en annexe, nous apporte une leçon nouvelle au vers 1047 : le `μοι' du \italique{Ravennas} n'y figure pas. De plus, il confirme -- comme l'avait déjà fait Willems en avance sur son temps -- le `βάρβαρον' du vers 1051 qui laissait de nombreux éditeurs sceptiques. (Voir sa \nieme{26} note, très instructive sur ce passage et le rôle de l'étoile Sirius.)}
